% Generated draft: requires mathematical and source review.
\documentclass[11pt]{article}
\usepackage[margin=25mm]{geometry}
\usepackage{fontspec}
\setmainfont{Latin Modern Roman}
\usepackage{amsmath,amssymb,mathtools,unicode-math}
\setmathfont{Latin Modern Math}
\usepackage{enumitem,needspace,xurl,hyperref}
\setlength{\parindent}{0pt}
\setlength{\parskip}{5pt}
\setlength{\emergencystretch}{3em}
\newcommand{\N}{\mathbb{N}}
\newcommand{\Z}{\mathbb{Z}}
\newcommand{\Q}{\mathbb{Q}}
\newcommand{\R}{\mathbb{R}}
\newcommand{\C}{\mathbb{C}}
\newcommand{\F}{\mathbb{F}}
\newcommand{\A}{\mathbb{A}}
\newcommand{\PP}{\mathbb{P}}
\newcommand{\E}{\mathbb{E}}
\newcommand{\eps}{\varepsilon}
\newcommand{\dd}{\,\mathrm{d}}
\begin{document}
\section*{UnsolvedMath ID: 30002406}
\textbf{Source problem code:} catalogue entry
\subsection{Definitions and mathematical statement}
For an integer \(d\geq 1\), write \(\mathbb{T}^{d}=\mathbb{R}^{d}/\mathbb{Z}^{d}\). To express the question conditionally, let \(\Gamma\) be a specified group or semigroup and let \(\alpha=(\alpha^\gamma)_{\gamma\in\Gamma}\) be a specified measurable action on \(\mathbb{T}^{d}\). In the invariant-probability-measure formulation, define
\[
\mathcal{M}_{\alpha}=\left\{\mu:\mu\text{ is a Borel probability measure on }\mathbb{T}^{d},\quad (\alpha^\gamma)_*\mu=\mu\text{ for every }\gamma\in\Gamma\right\}.
\]
Here invariance means that \(\mu((\alpha^\gamma)^{-1}(B))=\mu(B)\) for every Borel set \(B\subseteq\mathbb{T}^{d}\) and every \(\gamma\in\Gamma\). Probability normalization is conventional in this formulation but is not explicitly imposed by the supplied record.

\textbf{Source question.} For higher-rank toral actions, including the Furstenberg--Rudolph setting, obtain a partial classification of invariant measures under an explicit hypothesis replacing positive entropy, with particular attention to zero-entropy measures. The normalized statement asks for a hypothesis weaker than positive entropy, whereas the original statement asks for a replacement assumption without specifying a logical implication.

For \(\mu\in\mathcal{M}_{\alpha}\), let \(h_\mu(\alpha^\gamma)\) denote the Kolmogorov--Sinai entropy of the measure-preserving transformation \(\alpha^\gamma\). If zero entropy is intended elementwise, the corresponding class would be
\[
\mathcal{M}_{\alpha,0}=\left\{\mu\in\mathcal{M}_{\alpha}:h_\mu(\alpha^\gamma)=0\text{ for every }\gamma\in\Gamma\right\}.
\]
This is a conditional interpretation, not a hypothesis fixed by the source. The record does not specify whether entropy is measured for individual elements, selected generators, or the action itself.

The requested classification is a research objective rather than a fully specified conjecture. The source does not fix the acting group or semigroup, the meaning of higher rank, the permitted toral maps, ergodicity assumptions, the replacement hypothesis, or the desired classification conclusion. It also does not explicitly define the Furstenberg--Rudolph setting. These choices remain unresolved; the displayed framework does not assert a classification for every action introduced above.
\subsection{Short English statement}
The question asks whether invariant measures of higher-rank toral actions, including the Furstenberg--Rudolph setting, admit a partial classification beyond the positive-entropy regime, especially at zero entropy. The supplied record leaves the action class, entropy convention, replacement hypothesis, and classification target unspecified, so it does not determine a single precise conjecture.
\subsection{Sources}
\begin{itemize}
\item Local registry: \texttt{data/unsolvedmath/problems.json}, with \texttt{catalog/problems.json} designated for additional existing records; supplied permanent UnsolvedMath ID \(30002406\), exact source problem code \texttt{catalogue entry}, title ``Zero Entropy Measures for Higher Rank Toral Actions.'' Registry files were not independently inspected for this draft.
\item Supplied report citation: \emph{Noncommutative Geometry (2014)}, Oberwolfach Reports, DOI \texttt{10.4171/owr/2013/45}, pp.\ 2553--2629. \url{https://doi.org/10.4171/owr/2013/45}. The record describes the report as documenting workshop discussions and open problems. It supplies proposed year \(2013\), but no named proposer.
\item Supplied literature source, ``Oberwolfach Report'': \url{https://ems.press/content/serial-article-files/46476}.
\item Supplied literature source, ``Rudolph'': \url{https://doi.org/10.1017/S0143385700005629}.
\item Supplied literature source, ``Einsiedler survey'': \url{https://people.math.ethz.ch/~einsiedl/E-ICM2010.pdf}.
\item Literature assessment dated exactly \texttt{2026-08-21}: the supplied curation labels the problem partially solved and states that Rudolph's theorem gives the principal positive-entropy classification, while the requested zero-entropy or broader generalizations remain unresolved. This is supplied dated curation, not an independently verified literature assessment. The record explicitly requests independent verification of time-sensitive status and cited claims.
\item The registry labels the statement \texttt{corrected\_verified} and says that mathematical content was retained while notation, typography, and wording were normalized. This is supplied verification metadata, not independent mathematical verification by this draft.
\end{itemize}
\end{document}
